\section*{第 \ref{ch:b1:work-and-energy} 章——功、能量与势能}

\begin{solution}{exo:b1:work-and-energy:1}
提升力 $+mgh = 20 \times 9.81 \times 3.0 = \qty{589}{J}$；重力
$-\qty{589}{J}$；功率 $589/5.0 = \qty{118}{W}$。摩擦力：$-\mu_dmgd =
-0.30 \times 196 \times 3.0 = \qty{-177}{J}$。
\end{solution}

\begin{solution}{exo:b1:work-and-energy:2}
$E_k = \tfrac12 \times 1000 \times 27.8^2 = \qty{386}{kJ}$；$d = E_k/F =
\qty{55}{m}$；$v$ 加倍时 $E_k$ 和 $d$ 都变为四倍：\qty{220}{m}。
\end{solution}

\begin{solution}{exo:b1:work-and-energy:3}
$E_p = \tfrac12 \times 500 \times 0.10^2 = \qty{2.5}{J}$；$v = \sqrt{2 \times
2.5/0.050} = \qty{10}{m/s}$。
\end{solution}

\begin{solution}{exo:b1:work-and-energy:4}
$\tfrac12 mv^2 = mg\ell(1 - \cos 60^\circ)$：$v = \sqrt{g\ell} = \qty{3.1}{m/s}$。
张力：$T - mg = mv^2/\ell = mg$，$T = 2mg$。
\end{solution}

\begin{solution}{exo:b1:work-and-energy:5}
摩擦力做功 $-\mu_dmg\cos\alpha\,(h/\sin\alpha)$，故
\[
\tfrac12 mv^2 = mgh(1 - \mu_d\cot\alpha) = mgh(1 - 0.35),
\]
$v = \sqrt{2 \times 9.81 \times 2.0 \times 0.65} = \qty{5.1}{m/s}$；损失 $35\%$。
\end{solution}

\begin{solution}{exo:b1:work-and-energy:6}
$E_m = \tfrac12 mv^2 - GMm/R \geq 0$ 时物体可到达无穷远
（此时 $E_p = 0$ 且 $v \to 0$）：$v_e = \sqrt{2GM/R}$，即
\[
v_e = \sqrt{\frac{2 \times 6.67\times10^{-11} \times 5.97\times10^{24}}{6.37\times10^6}}
= \qty{11.2}{km/s}.
\]
密度相同，$M \propto R^3$：$v_e \propto R$，故 \qty{22.4}{km/s}。
\end{solution}

\begin{solution}{exo:b1:work-and-energy:7}
最低点：$mg(\ell_0 + x) = \tfrac12 kx^2$，$25x^2 - 687x - 13734 = 0$，
$x = \qty{41}{m}$；总下落 \qty{61}{m}；$T_{\max} = kx = \qty{2.0}{kN}
= 3.0\,mg$（合力 \qty{2}{g} 向上）。速率在合力
为零处最大，即 $kx = mg$：$x = \qty{14}{m}$，位于跳台下方 \qty{34}{m} 处。
\end{solution}

\begin{solution}{exo:b1:work-and-energy:8}
重力：$mgv\sin\alpha = 80 \times 9.81 \times 5.0 \times 0.05 = \qty{196}{W}$；
阻力：$\tfrac12\rho C_xSv^3 = 0.5 \times 1.2 \times 0.50 \times 125 =
\qty{38}{W}$；总计 \qty{234}{W}。
\end{solution}

\begin{solution}{exo:b1:work-and-energy:9}
$E_p' = 4E_0(x^3/a^4 - x/a^2) = 0$：$x = 0, \pm a$。$E_p'' = E_0(12x^2/a^4 -
4/a^2)$：$-4E_0/a^2$ 在 $0$ 处（不稳定极大值，$E_p = 0$），$+8E_0/a^2$
在 $\pm a$ 处（稳定极小值，$E_p = -E_0$）：深度为 $E_0$ 的\omterm{def:b1:work-and-energy:turning}{势阱}；
$\omega_0^2 = 8E_0/ma^2$。
\end{solution}

\begin{solution}{exo:b1:work-and-energy:10}
$\tfrac12 m\dot x^2 + \tfrac12 kx^2 = E$：椭圆，半轴为 $x_{\max}
= \sqrt{2E/k}$ 与 $\dot x_{\max} = \sqrt{2E/m}$。运动为\omterm{prop:b1:point-kinematics:harmonic}{简谐运动}，
周期为 $2\pi\sqrt{m/k}$，与\omterm{def:b1:wave-propagation:signal}{振幅}无关（等时性），故
各椭圆所用时间相同。
\end{solution}

\begin{solution}{exo:b1:work-and-energy:11}
能量：$v^2 = 2gR(1 - \cos\theta)$。径向：$mg\cos\theta - N = mv^2/R$，
故 $N = mg(3\cos\theta - 2)$，在 $\cos\theta = 2/3$（$\theta =
48^\circ$）处为零：高度 $2R/3$，速率 $\sqrt{2gR/3}$。
\end{solution}

\begin{solution}{exo:b1:work-and-energy:12}
$\tfrac12 mv^2 = \mu_dmgd$：$\mu_d = 625/(2 \times 9.81 \times 80) = 0.40$。
轮胎与路面接触处产生热量（并留下刹车痕迹）。滚动的车轮受
静摩擦力作用，$\mu_s > \mu_d$：制动力更大，
制动距离更短，并保持转向能力。
\end{solution}

\begin{solution}{pb:b1:work-and-energy:1}
\textbf{1.} $r \to 0$：$+\infty$（$r^{-12}$ 项）；$r \to \infty$：
$0^-$。陡峭的壁、一个\omterm{def:b1:work-and-energy:turning}{势阱}、一条长尾。

\textbf{2.} $E_p' = \epsilon[-12\sigma^{12}/r^{13} + 12\sigma^6/r^7] = 0$ 当且仅当
$(\sigma/r)^6 = 1$：$r = \sigma$，$E_p = \epsilon(1 - 2) = -\epsilon$。

\textbf{3.} $(\sigma/r)^{12} = u^2$：$E_p = \epsilon(u^2 - 2u) = \epsilon[(u - 1)^2
- 1]$，在 $u = 1$（$r = \sigma$）处取最小值 $-\epsilon$。

\textbf{4.} $F = -E_p' = (12\epsilon/\sigma)[(\sigma/r)^{13} - (\sigma/r)^7]$：
$r < \sigma$ 时为正（斥力），更远处为负（引力）。

\textbf{5.} $(1/1.2)^6 = 0.335$，$(1/1.2)^{12} = 0.112$：$E_p = \epsilon(0.112
- 0.670) = -0.56\epsilon = \qty{-2.5}{eV}$。$(1/1.2)^{13} = 0.0935$，
$(1/1.2)^7 = 0.279$：$F = (12\epsilon/\sigma)(-0.186) = -2.2\epsilon/\sigma =
-2.2 \times 7.0\times10^{-19}/1.27\times10^{-10} = \qty{-12}{nN}$。

\textbf{6.} $\epsilon = \qty{4.4}{eV} = 4.4 \times 96.5 = \qty{425}{kJ/mol}$，
接近实测值 \qty{431}{kJ/mol}。

\textbf{7.} $E_p'' = \epsilon[156\sigma^{12}/r^{14} - 84\sigma^6/r^8]$；在
$\sigma$ 处：$72\epsilon/\sigma^2 > 0$：稳定的。

\textbf{8.} $E_p \approx -\epsilon + \tfrac12(72\epsilon/\sigma^2)(r - \sigma)^2$：
$k = 72\epsilon/\sigma^2$。

\textbf{9.} $k = 72 \times 7.0\times10^{-19}/(1.27\times10^{-10})^2 =
\qty{3.1e3}{N/m}$。

\textbf{10.} $\omega_0 = \sqrt{k/m} = \sqrt{3.1\times10^3/1.67\times10^{-27}} =
\qty{1.4e15}{rad/s}$; $f = \qty{2.2e14}{Hz}$; $\lambda = c/f = \qty{1.4}{\micro m}$;
$1/\lambda = \qty{7300}{cm^{-1}}$.

\textbf{11.} 偏高 $2.4$ 倍：模型的壁过陡
（$r^{-12}$ 芯部只是共价键的粗略替代），但
\omterm{def:b1:units-dimensions:oom}{数量级}——$10^{14}$ Hz 的红外振动——
是对的，而且所用方法（\omterm{def:b1:work-and-energy:turning}{势阱}的曲率）与
采用更好势函数时完全相同。

\textbf{12.} $\tfrac12 kA^2 \approx k_BT$: $A = \sqrt{2k_BT/k} = \sqrt{8.3\times
10^{-21}/3.1\times10^3} = \qty{1.6e-12}{m} = 0.013\sigma$：键
几乎只是微微颤动。

\textbf{13.} $v_{\max} = A\omega_0 = 1.6\times10^{-12} \times 1.4\times10^{15}
= \qty{2.2}{km/s}$.

\textbf{14.} $-\epsilon < E_m < 0$：围绕 $(\sigma, 0)$ 的闭曲线——
振动，束缚；$E_m > 0$：开曲线——原子分离；
$E_m = 0$：分界线，原子恰好能够到达无穷远。

\textbf{15.} $u^2 - 2u = -\tfrac12$，其中 $u = (\sigma/r)^6$：$u = 1 \pm 1/\sqrt2
= 1.71$ 或 $0.293$；$r = \sigma u^{-1/6}$：$0.915\sigma$ 和 $1.23\sigma$。

\textbf{16.} 中点 $1.07\sigma > \sigma$：外侧势壁比内侧更软，
因此原子在外侧活动的范围更大；时间平均的
距离随能量增大——受热的键变长，固体膨胀。

\textbf{17.} 周期更长：外侧较软一侧的回复力
比谐振子近似所假定的更弱，因此回到平衡位置
所用时间更长。

\textbf{18.} $\epsilon/2 = \qty{2.2}{eV} = \qty{3.5e-19}{J}$；作为
$\sigma$ 处的\omterm{thm:b1:work-and-energy:ke}{动能}：$v = \sqrt{2 \times 3.5\times10^{-19}/1.67\times10^{-27}}
= \qty{2.0e4}{m/s}$。

\textbf{19.} 能量守恒：$E_m > 0$ 保持 $> 0$，曲线为
开曲线，两原子一次碰撞后即飞散。第三个物体（另一个
分子、一堵壁或一个光子）必须带走多余的能量。

\textbf{20.} $W = E_p(2\sigma) - E_p(\sigma) = \epsilon(2^{-12} - 2 \times 2^{-6})
+ \epsilon = 0.97\epsilon = \qty{4.3}{eV}$，为正：吸引力将
使原子向内靠近，是动力。

\textbf{21.} $W_{\text{外力}} = E_p(1.5\sigma) - E_p(\sigma) = \epsilon(0.0077 -
0.176) + \epsilon = 0.83\epsilon = \qty{3.7}{eV}$。

\textbf{22.} $(1/0.9)^{13} = 3.93$，$(1/0.9)^7 = 2.09$：$F = +22\epsilon/\sigma
= \qty{+120}{nN}$，是 $1.2\sigma$ 处吸引力的十倍：压缩一个
键 $10\%$ 所付出的代价远大于拉伸它 $20\%$。

\textbf{23.} $T = \epsilon/k_B = 7.0\times10^{-19}/1.38\times10^{-23} =
\qty{5.1e4}{K}$。真实气体的解离温度远低于此值，因为能量
分布不均：一部分分子携带的能量是 $k_BT$ 的许多倍
（见\cref{ch:b1:kinetic-theory}）。

\textbf{24.} 保持不变的：最小值处的平衡，$k = E_p''$ 于
最小值处，解离能 $=$ \omterm{def:b1:work-and-energy:turning}{势阱}深度，\omterm{def:b1:work-and-energy:turning}{转折点}由 $E_p = E_m$ 给出，
不对称 $\Rightarrow$ 膨胀。改变的是：$k$ 的数值、
频率和\omterm{def:b1:work-and-energy:turning}{转折点}的数值。

\textbf{25.} 平衡：$E_p' = 0$；振动：$\omega_0^2 = E_p''/m$；
逃逸：将 $E_m$ 与\omterm{def:b1:work-and-energy:turning}{势阱}边缘相比较。
\end{solution}
